\chapter{Discussion}
\label{ch:discussion}

This chapter interprets the experimental results of
Chapter~\ref{ch:experiments} in the light of the research questions, examines
the threats to the validity of the conclusions, and delimits the scope within
which the findings can be generalised.

\promoterinline{This chapter is a scaffold. The interpretive claims below are
written as guiding prompts and are to be completed once the full 600-image run
populates the tables and figures of Chapter~\ref{ch:experiments}. No empirical
claim should be asserted here until the corresponding result cell is filled.}

%% ========================================================================
\section{Interpretation of the Rankings}
\label{sec:disc-rankings}

This section will discuss what the aggregated rankings of
Section~\ref{sec:exp-ranking} imply for practitioners selecting an FAE method.

\begin{itemize}
    \item \textbf{Does multi-criteria aggregation change the verdict?}
    Compare the FAE-method ordering under the single-metric baseline against the
    MetaQuantus-discounted index (Table~\ref{tab:rank-stability}). Discuss the
    cases in which a method that ``wins'' on Faithfulness alone is demoted once
    Robustness, Localization, and Randomization enter the aggregate.
    \item \textbf{Effect of the weighting scheme.} Contrast uniform,
    Autoweighted, and MQ-discounted weights. Where the three agree, argue that
    the conclusion is robust to the aggregation choice; where they disagree,
    identify which metrics drive the divergence and whether the divergence
    tracks the reliability discount.
    \item \textbf{Per-category trade-offs.} Use the ensembling comparison
    (Section~\ref{sec:exp-ensemble}) to discuss whether NormEnsembleXAI buys
    Robustness at the expense of Localization or Complexity, and whether the
    scalar index hides such trade-offs.
    \item \textbf{Cross-model consistency.} Discuss whether the ResNet-18 and
    SqueezeNet rankings agree; persistent agreement strengthens the claim that
    the framework measures a method-intrinsic property rather than a
    model-specific artefact.
\end{itemize}

%% ========================================================================
\section{Implications of the Redundancy and Reliability Findings}
\label{sec:disc-redundancy}

\begin{itemize}
    \item \textbf{Redundant categories.} The pilot indicates near-perfect
    correlation within the Complexity and Robustness categories and near-zero
    correlation within Faithfulness and Localization. Discuss the implication:
    reporting both members of a redundant pair overstates evidence, and the
    orthogonal categories carry the discriminative signal.
    \item \textbf{Greedy cross-category pruning.} The current pruning is greedy
    and can drop a metric because it correlates with a metric in a
    \emph{different} category (Outstanding Issue~2 in the redundancy report).
    Discuss whether within-category pruning (Decision~D4 as originally stated)
    would yield a more defensible $M^{*}$ and how the choice affects the
    downstream ranking.
    \item \textbf{Reliability vs.\ agreement.} Discuss whether the
    MetaQuantus reliability ordering coincides with the ensemble-agreement
    ordering. If a metric is highly agreeing but unreliable (or vice versa), the
    MQ-discount changes its weight; quantify how often this occurs.
\end{itemize}

%% ========================================================================
\section{Threats to Validity}
\label{sec:disc-threats}

\subsection{Internal Validity}
\begin{itemize}
    \item \textbf{Sample size.} The redundancy and reliability findings reported
    in the pilot rest on 12 images. Correlation estimates and the resulting
    pruning decisions are unstable at this $n$; the membership of $M^{*}$ may
    change on the full run. This is the single largest threat and is the reason
    every pilot value is labelled as such.
    \item \textbf{NaN handling in MPRT.} Model Parameter Randomisation produced
    many NaNs for SqueezeNet (only 27.4\% valid observations in the pilot) and
    for Grad-CAM on randomised ResNet. Discuss whether the surviving
    observations are representative or a biased subsample, and whether the
    differing MPR/sparseness correlation between the two models is a true
    architectural effect or a sparse-data artefact.
    \item \textbf{Target-class choice.} Attributions target the model's
    prediction, not the ground-truth label; discuss how this interacts with
    misclassified images.
\end{itemize}

\subsection{Construct Validity}
\begin{itemize}
    \item \textbf{Are the metrics measuring what they claim?} The
    meta-evaluation partially addresses this, but NR and AR are themselves
    proxies. Discuss the residual risk that a metric passes meta-evaluation yet
    fails to capture human-meaningful explanation quality.
    \item \textbf{Normalisation choice.} Second Moment Scaling (Decision~D1) was
    adopted from NormEnsembleXAI; discuss sensitivity of the ranking to the
    alternative $z$-score and min--max normalisations implemented for ablation.
\end{itemize}

\subsection{External Validity}
\begin{itemize}
    \item \textbf{Single dataset and domain.} All evidence comes from ISIC~2017
    dermoscopy. Discuss the extent to which conclusions transfer to other
    medical imaging modalities or to natural images, and why the medical domain
    was nonetheless chosen (mask availability, motivation).
    \item \textbf{Two architectures.} ResNet-18 and SqueezeNet are both
    convolutional; transformer-based classifiers may interact differently with
    gradient-based attributions and with MPRT.
\end{itemize}

%% ========================================================================
\section{Limitations}
\label{sec:disc-limitations}

\begin{itemize}
    \item Computational cost of the perturbation-based metrics and of LIME/
    KernelSHAP constrained the method panel; discuss what was excluded and why.
    \item The framework evaluates attribution \emph{maps}; it does not address
    concept-level or counterfactual explanations.
    \item Human-grounded validation (clinician studies) is explicitly out of
    scope; the effectiveness index is functionally grounded only.
\end{itemize}
